\chapter{Orientación y elección de ruta}
\label{chap:orientacion}

\section{Qué componentes intervienen}

La AKD1000 no se maneja como una aplicación aislada. El host ejecuta el sistema operativo y el
runtime; PCIe transporta órdenes y datos; el controlador expone el dispositivo al sistema; y el
modelo FBZ define la red que el runtime puede mapear. Separar estas capas evita buscar un fallo
de alimentación en Python o intentar resolver un paquete ausente reinstalando el controlador.

\begin{table}[H]
\centering
\caption{Responsabilidad de cada capa del sistema.}
\label{tab:capas}
\begin{tabularx}{\textwidth}{@{}L{0.22\textwidth}YY@{}}
\toprule
\textbf{Capa} & \textbf{Qué hace} & \textbf{Qué no demuestra por sí sola} \\
\midrule
Host & Ejecuta Linux o Windows, Python, preparación de datos y registros. & Que el cálculo se haya realizado en la tarjeta. \\
Interfaz PCIe & Transporta configuración, entradas y salidas. & Que exista un controlador utilizable. \\
Controlador Linux & Crea el acceso al dispositivo, normalmente mediante \code{/dev/akida*}. & Que el modelo sea compatible. \\
Runtime Akida & Carga FBZ, descubre dispositivos, mapea y ejecuta. & Que un mapeo virtual sea una inferencia física. \\
AKD1000 & Ejecuta las secuencias asignadas al hardware. & El tiempo total de Python ni el consumo de toda la plataforma. \\
FBZ & Conserva una red Akida serializada. & El significado de las entradas, las etiquetas o el dataset. \\
\bottomrule
\end{tabularx}
\end{table}

La tarjeta PCIe documentada por BrainChip contiene un AKD1000, usa una interfaz host PCIe y
dispone de memoria LPDDR4 accesible mediante DMA \cite{brainchip_card_brief}. El product brief
del SoC describe veinte núcleos de procesamiento neuronal y un endpoint PCIe 2.1 de dos líneas
\cite{brainchip_soc_brief}. Estas son características del producto; no son una promesa de
latencia, energía o utilización para cualquier red.

\section{Tres decisiones antes de empezar}

\paragraph{Decisión 1: software o hardware.}
Si aún no tienes una tarjeta, sigue la ruta \textbf{S}: entorno software, modelo de prueba y
dispositivo virtual. Si tienes la tarjeta y una conexión PCIe documentada, completa primero la
ruta S y continúa con la ruta \textbf{H}. La simulación confirma el modelo y el paquete; no
confirma el cableado, el controlador ni la placa.


\paragraph{Decisión 2: entorno actual o de referencia.}
Para aprender con la documentación vigente usa la ruta \textbf{A}: MetaTF 2.19.3, o la versión
que figure en la página oficial cuando leas el manual. Para repetir la prueba incluida usa la ruta
\textbf{R}: Python 3.11 y Akida 2.19.1.
No mezcles ambas instalaciones
en el mismo entorno virtual.


\paragraph{Decisión 3: tipo de host.}
Un PC Linux con ranura PCIe y una Raspberry Pi 5 con adaptador FFC requieren montajes distintos.
Windows está incluido en las configuraciones software actuales de MetaTF
\cite{brainchip_install_2193}, pero este manual no afirma soporte de la tarjeta física en
Windows: las instrucciones del controlador PCIe público consultado son para Linux
\cite{brainchip_driver}.


\section{Matriz de rutas}

\begin{longtable}{@{}L{0.12\textwidth}L{0.21\textwidth}L{0.29\textwidth}L{0.28\textwidth}@{}}
\caption{Recorridos cubiertos por el manual.}\label{tab:rutas}\\
\toprule
\textbf{Ruta} & \textbf{Equipo} & \textbf{Objetivo} & \textbf{Evidencia final} \\
\midrule
\endfirsthead
\toprule
\textbf{Ruta} & \textbf{Equipo} & \textbf{Objetivo} & \textbf{Evidencia final} \\
\midrule
\endhead
S & Windows o Linux compatible, sin tarjeta & Instalar MetaTF, importar Akida y ejecutar en software. & Versión registrada, salida del modelo y mapeo virtual. \\
H-PC & Host Linux, tarjeta y conexión PCIe documentada & Detectar y ejecutar en AKD1000 mediante el controlador público. & \code{lspci}, módulo, nodo, \code{akida devices} y prueba física. \\
H-Pi & Raspberry Pi 5, adaptador FFC y tarjeta & Habilitar PCIe cuando proceda y completar H-PC. & Configuración de arranque, enumeración y prueba física. \\
R & Python 3.11 y SDK de referencia & Repetir la prueba básica con Akida 2.19.1. & Salida conocida y mapeo virtual comprobado. \\
\bottomrule
\end{longtable}

\section{Comprobar cada etapa}

Cada etapa responde a una pregunta distinta:

\begin{enumerate}[nosep]
  \item \textbf{Inventario:} ¿sé exactamente qué host, tarjeta, adaptador, fuentes, sistema,
        kernel y arquitectura tengo?
  \item \textbf{Montaje:} ¿la cadena está conectada según la documentación de esas revisiones?
  \item \textbf{Enumeración:} ¿el host ve un endpoint PCIe nuevo?
  \item \textbf{Controlador:} ¿el módulo se ha compilado para este kernel y existe acceso al nodo?
  \item \textbf{SDK:} ¿el intérprete activo importa la versión prevista de Akida?
  \item \textbf{Modelo:} ¿el FBZ carga, tiene el contrato esperado y mapea como se pretende?
  \item \textbf{Ejecución:} ¿una entrada conocida produce una salida conocida en hardware?
  \item \textbf{Medida:} ¿la ventana, las unidades y la frontera de potencia están declaradas?
\end{enumerate}

Resuelve la primera etapa que falla antes de continuar. Guarda su error para localizar
el problema.
